% Options for packages loaded elsewhere
\PassOptionsToPackage{unicode}{hyperref}
\PassOptionsToPackage{hyphens}{url}
\PassOptionsToPackage{dvipsnames,svgnames,x11names}{xcolor}
%
\documentclass[
  11pt,
  a4paper,
]{ltjsarticle}
\usepackage{amsmath,amssymb}
\usepackage{iftex}
\ifPDFTeX
  \usepackage[T1]{fontenc}
  \usepackage[utf8]{inputenc}
  \usepackage{textcomp} % provide euro and other symbols
\else % if luatex or xetex
  \usepackage{unicode-math} % this also loads fontspec
  \defaultfontfeatures{Scale=MatchLowercase}
  \defaultfontfeatures[\rmfamily]{Ligatures=TeX,Scale=1}
\fi
\usepackage{lmodern}
\ifPDFTeX\else
  % xetex/luatex font selection
\fi
% Use upquote if available, for straight quotes in verbatim environments
\IfFileExists{upquote.sty}{\usepackage{upquote}}{}
\IfFileExists{microtype.sty}{% use microtype if available
  \usepackage[]{microtype}
  \UseMicrotypeSet[protrusion]{basicmath} % disable protrusion for tt fonts
}{}
\makeatletter
\@ifundefined{KOMAClassName}{% if non-KOMA class
  \IfFileExists{parskip.sty}{%
    \usepackage{parskip}
  }{% else
    \setlength{\parindent}{0pt}
    \setlength{\parskip}{6pt plus 2pt minus 1pt}}
}{% if KOMA class
  \KOMAoptions{parskip=half}}
\makeatother
\usepackage{xcolor}
\usepackage[margin=24mm]{geometry}
\setlength{\emergencystretch}{3em} % prevent overfull lines
\providecommand{\tightlist}{%
  \setlength{\itemsep}{0pt}\setlength{\parskip}{0pt}}
\setcounter{secnumdepth}{5}
\ifLuaTeX
  \usepackage{selnolig}  % disable illegal ligatures
\fi
\IfFileExists{bookmark.sty}{\usepackage{bookmark}}{\usepackage{hyperref}}
\IfFileExists{xurl.sty}{\usepackage{xurl}}{} % add URL line breaks if available
\urlstyle{same}
\hypersetup{
  colorlinks=true,
  linkcolor={blue},
  filecolor={Maroon},
  citecolor={Blue},
  urlcolor={blue},
  pdfcreator={LaTeX via pandoc}}

\author{}
\date{}

\begin{document}

\hypertarget{ux90e8ux5206ux7a7aux9593}{%
\section{部分空間}\label{ux90e8ux5206ux7a7aux9593}}

\texttt{V} の部分集合 \texttt{W} が部分空間であるとは、\texttt{W}
自身がベクトル空間になっていることです。

実用上は、

\begin{enumerate}
\def\labelenumi{\arabic{enumi}.}
\tightlist
\item
  \texttt{0∈W}
\item
  \texttt{x,y∈W\ ⇒\ x+y∈W}
\item
  \texttt{x∈W,\ α∈R\ ⇒\ αx∈W}
\end{enumerate}

を確認します。

\hypertarget{ux7ddaux5f62ux4ee3ux6570ux3067ux983bux51faux3059ux308bux90e8ux5206ux7a7aux9593}{%
\subsection{線形代数で頻出する部分空間}\label{ux7ddaux5f62ux4ee3ux6570ux3067ux983bux51faux3059ux308bux90e8ux5206ux7a7aux9593}}

行列 \texttt{A} には重要な4つの部分空間があります。

\begin{itemize}
\tightlist
\item
  列空間 \texttt{Col(A)}
\item
  行空間 \texttt{Row(A)}
\item
  零空間 \texttt{Null(A)}
\item
  左零空間 \texttt{Null(A\^{}T)}
\end{itemize}

SVD はこれらを直交基底によって同時に整理します。

特に零空間は「変換によって完全に消える入力方向」、列空間は「変換後に到達可能な出力方向」です。この幾何学的解釈を常に意識してください。

\end{document}
